\documentclass[11pt,a4paper]{article}

\input{preamble.tex}

\title{%
  \AddToHook{shipout/background}{%
  \ifnum\value{page}=1
    \uitTitleLogo
  \fi
}
  \Large
  \textcolor{uitDarkBlue}{%
    \textbf{Writing Scientific Reports in\\
    Computer Science and Engineering}%
  }
  \\[0.3em]
  \large
  A Practical Guide and \LaTeX{} Template for Bachelor Students
}

\author{
  Course: \emph{DTE-2602 Introduction to Machine Learning and Artifical Intelligence}\\
  Instructor: \emph{Vegard Seines}
}

\date{\today}

\begin{document}

\maketitle

\begin{abstract}
  In this assignment Djikstras algorithm, a graph with a start and goal was given as well as a placeholder A* function. The objective was to implement the A* algorithm and compare the different solutions to a provided path finding problem. The A* algorithm was implementation was based on the provided Djikstra algorithm and tested with the provided program. The A* algorithm and Djikstra both found the same optimal path. Djikstra visited 196 nodes while A* only visited 163.
\end{abstract}

\newpage
\tableofcontents
\thispagestyle{fancy}
\newpage

% Its not always necessary to add a \newline after each section, 
% but it is often done. It makes it easier to see where each new section starts.
\section{Introduction}
\label{sec:introduction}

The given assignment comprised of multiple parts.
\begin{itemize}
  \item Implement the A* placeholder.
  \item Test that the A* implementation works on the provided graph.
  \item Run A* with the same start and goal nodes as well as the heuristikk that is used by Dijkstra.
  \item Compare the results of the A* and Djikstra run.
  \item Discuss how A* is different from Greedy Search.
\end{itemize}

\newpage
\section{Background}
\label{sec:background}

Path finding is an important area in computer science empowering all kinds of technology from GPS and navigation to network topology and more abstract optimization problems.
Djikstras algorithm and A* are two important algorithms for navigating weighted graphs.

\subsection{Weighted Graphs}

A graph is a collection of nodes and edges and a weighted graph has a cost associated with traversing the edges. This simple data structure is profoundly usefull for describing a host of impactfull use cases. There are several ways to represent a graph which offers different advantages in terms of memory layout and the efficency of the consuming algorithms.

\subsubsection{Adjecancy List}
An adjacency list maps every node to a collection of its neighboring nodes and edge weights. In python this is usually implemented using a dictionary where keys are nodes and values are lists of tuples.
This structure is good for representing sparse graphs where nodes have few connections and can be a space efficient way to store a graph. Querying neighbors for a specific node is a fast operation in adjecancy lists.

\subsection{Djikstras algorithm}

TODO

\subsection{Greedy Search}

TODO

\subsection{A*}

TODO

\newpage
\input{sections/03_implementation}
\newpage
\input{sections/04_results}
\newpage
\input{sections/05_discussion}
\newpage
\input{sections/06_conclusion}
\newpage

\addcontentsline{toc}{section}{References}
\bibliographystyle{ieeetr}
\bibliography{references}

\end{document}
